% STAGING ONLY.  This statement is not part of the live manuscript.
% It must not be integrated before exact-byte production Rethlas verification.
\begin{theorem}\label{thm:staging-basepoint-free-relative-resultant}
Let \(S\) be a locally Noetherian scheme and let
\(f\colon X\to S\) be a flat projective morphism whose fibres are nonempty
and purely \(n\)-dimensional.  For \(0\leq i\leq n\), let \(L_i\) be an
\(f\)-ample line bundle, let \(E_i\) be a finite locally free
\(\mathcal O_S\)-module of positive rank, and fix a surjection
\begin{equation}\label{eq:staging-resultant-evaluation}
 f^*E_i\twoheadrightarrow L_i.
\end{equation}
Let
\[
 P_i:=\mathbf P_S(E_i^\vee),
 \qquad q_i\colon P_i\to S.
\]
We use the convention that \(P_i\) carries the tautological subline
\(\mathcal O_{P_i}(-1)\hookrightarrow q_i^*E_i\).  Put
\[
 P:=\prod_{i=0}^n{}_S P_i,
 \qquad q\colon P\to S,
 \qquad Y:=X\times_S P,
\]
and denote by \(p_i\colon P\to P_i\), \(g\colon Y\to X\), and
\(\pi\colon Y\to P\) the projections.  Set
\begin{equation}\label{eq:staging-resultant-universal-line}
 \widetilde L_i
 :=g^*L_i\otimes\pi^*p_i^*\mathcal O_{P_i}(1).
\end{equation}
The tautological subline and \eqref{eq:staging-resultant-evaluation}
define a universal section
\(\sigma_i\in H^0(Y,\widetilde L_i)\).

Let \(I_{X/S}\) denote Ducrot's symmetric, multiadditive intersection-line
functor.  There is a canonical section
\begin{equation}\label{eq:staging-resultant-section}
 \operatorname{Res}_{E_\bullet}(\sigma_0,\ldots,\sigma_n)
 \in H^0(P,\mathscr R),
 \qquad
 \mathscr R:=I_{Y/P}(\widetilde L_0,\ldots,\widetilde L_n),
\end{equation}
with the following properties.

For \(0\leq i\leq n\), let
\begin{equation}\label{eq:staging-resultant-mixed-degree}
 d_i(s):=
 \bigl(L_{0,s}\cdots\widehat{L_{i,s}}\cdots L_{n,s}\bigr)_{X_s}.
\end{equation}
This is a locally constant positive integer-valued function on \(S\).
Writing \(\mathcal O_{P_i}(d_i)\) componentwise on the open-and-closed
rank-and-degree strata of \(S\), there is a canonical isomorphism
\begin{equation}\label{eq:staging-resultant-line-formula}
 \mathscr R\simeq
 q^*I_{X/S}(L_0,\ldots,L_n)
 \otimes
 \bigotimes_{i=0}^n p_i^*\mathcal O_{P_i}(d_i).
\end{equation}

For a permutation \(\tau\) of \(\{0,\ldots,n\}\), let \(U_\tau\subset P\)
be the open locus on which the ordered sequence
\[
 \bigl(\sigma_{\tau(0)},\ldots,\sigma_{\tau(n-1)}\bigr)
\]
is relative regular on \(Y/P\); when \(n=0\), the empty sequence is
understood and
\(U_\tau=P\).  On \(U_\tau\), successive restriction identifies
\(\mathscr R\) with the norm line of
\(\widetilde L_{\tau(n)}\) on the finite flat scheme
\[
 V(\sigma_{\tau(0)})\cap\cdots\cap
 V(\sigma_{\tau(n-1)})\longrightarrow U_\tau,
\]
and \eqref{eq:staging-resultant-section} is identified with the norm of
\(\sigma_{\tau(n)}\).  These local descriptions characterize the section
uniquely.

Let
\begin{equation}\label{eq:staging-resultant-incidence}
 \mathcal Z:=\pi\left(
   \bigcap_{i=0}^n V(\sigma_i)
 \right)\subset P.
\end{equation}
The image is closed.  If \(p\in P\), \(s=q(p)\), and
\begin{equation}\label{eq:staging-resultant-depth-firewall}
 \operatorname{depth}\mathcal O_{P_s,p}\leq1,
\end{equation}
then
\begin{equation}\label{eq:staging-resultant-incidence-detection}
 p\in\mathcal Z
 \quad\Longleftrightarrow\quad
 \operatorname{Res}_{E_\bullet}(\sigma_0,\ldots,\sigma_n)(p)=0.
\end{equation}
Equivalently, on every fibre the resultant detects the incidence locus at
all points of codimension at most one.  If \(n=0\),
\eqref{eq:staging-resultant-incidence-detection} holds for every \(p\in P\).

The line, its canonical section, the isomorphism
\eqref{eq:staging-resultant-line-formula}, and
\eqref{eq:staging-resultant-incidence-detection} are compatible with every
base change \(T\to S\) for which \(T\) is locally Noetherian.  They are also
compatible with permutations of the systems through the symmetry
isomorphisms of the intersection line.

Finally, on an open-and-closed component on which \(n=0\) and
\(\operatorname{rk}E_0=1\), the evaluation map is an isomorphism,
\(\widetilde L_0\simeq\mathcal O_Y\),
\(\mathscr R\simeq\mathcal O_P\), and the resultant is the unit section.
For \(n>0\), the hypotheses force \(\operatorname{rk}E_i\geq2\) for every
\(i\).
\end{theorem}
